\section*{Abstract}
This course report documents a questionnaire dataset collected from undergraduate engineering students in Bangladesh. Two independently administered Google Forms captured the same core academic, behavioural, lifestyle, motivation, support, and study-environment concepts with small differences in wording, language, ordering, and grade-band options. The source files contain 1,126 responses: 230 from BUET, BRAC University, and CUET, and 896 from KUET. A reproducible pipeline maps both forms to a common schema, removes source-specific text differences, normalises university and department names, reconciles grade bands, records missing-value decisions, and removes 18 records without a usable CGPA target. The resulting dataset contains 1,108 cleaned student records from four universities and represents both recent SGPA and cumulative CGPA as four ordered groups from C0 to C3. The project package includes raw and cleaned tabular files, WEKA-ready ARFF files, a data dictionary, cleaning logs, descriptive figures, fixed train/test indices, model-output tables, and executable notebooks. These resources support reuse in educational data analysis, questionnaire harmonisation, ordinal classification, interpretable modelling, and comparison of standard machine-learning algorithms. The latter part of this report demonstrates machine-learning analyses that can be performed on the dataset while keeping data description, predictive evaluation, and interpretation as separate tasks.

\textbf{Keywords:} academic performance; questionnaire data; engineering students; educational analytics; ordinal classification; interpretable machine learning; university survey

